\cotefile{12}
\scanpage{22}
% archive sheet number: 21
% Several mathematical slots in this typescript were never filled: spacing below preserves them without supplying symbols.
% A long diagonal cancellation runs from the third text line to the end of the first paragraph.
Or par raison de platitude $H'_i=\underline{\mathrm{Hom}}\,(\underline{O}_{X'}/J',G'_i)$, donc on peut \unsure{supposer} (quitte à remplacer les $G_i$ par les $H_i$) que les $G_i$ sont annulés par $J$. \deleted{Remplaçant maintenant $X$ par $\operatorname{Spec}(\underline{O}_X/\underline{J})$, on est ramené au cas où $F$ est un Module fidèle, donc $X$ est irréductible de point générique l'unique point $x\quad\operatorname{Ass}F$. En vertu de théorème de platitude générique, il existe un voisinage ouvert \unsure{uffine|affine|Uffine} $U$ de $x$ tel que l'anneau de son image inverse $U'$ dans $X'$ soit \uline{libre} sur celui de $U$. Passant \deleted{\unsure{au|en}} à l'interprétation en termes de modules sur ces anneaux, on voit alors aussitôt que \deleted{sur $U'$, l'intersection des} $\quad{}'$ est rulle sur $U'$. Or en vertu de \unsure{par} \handwritten{3.3.} on a $\operatorname{Ass}F'\quad X'_x\quad U'$, donc $\quad{}'=0$, Supposons maintenant $n\quad 1$, soit $x$ \inserted{un élément maximal de} $\operatorname{Ass}F$, d'après ce qui précède, \deleted{il existe un voisinag} si $U$ est un voisinage ouvert de $x$ ne contenant aucun autre \unsure{Élément|élément} de $\operatorname{Ass}F$, $U'$ son image inverse dans $U$, alors on a $\quad{}'\ U'=0$. En d'autres termes, si $i\colon U\quad X$ est l'imm can, et $v\colon F\quad i\ (F\ U)$ l'homomorphisme canonique, donc le transformé par $\otimes_X X'$ s'identifie à $v'\colon F'\quad i'\ (F'\ U')$, \deleted{\unsure{alors}} (réf), alors $\quad{}'$ est une section de $\operatorname{Ker}v'=\operatorname{Ker}v\otimes_X X'$ (re-platitude). Introduisant $F_1=\operatorname{Ker}v$, on peut donc considérer $\quad{}'$ comme une section de $F'_1$, et on est ramené au problème précédent avec $F$ remplacé par $F_1$. Mais $\operatorname{Ass}F_1=\operatorname{Ass}F-x$, et on gagne par hyp de récurrence. (NB au lieu de récurrer sur $\operatorname{card}\operatorname{Ass}$ on peut aussi dans ce dernier argument envisager qu'on a fait une récurrence \deleted{sur} noethérienne sur $\operatorname{supp}F$, ce qui est plutôt plus naturel et dispense de vérifier une formule à la noix sur $\operatorname{Ass}F_1$).}

Bien entendu, prop. \handwritten{3} a) (ii) résulte de là, car on y peut supposer $X$ et $X'$ affines, les $Z_i$ affines (quitte à remplacer chaque $Z_i$ par des ouverts affines le recouvrant), et on applique le critère 1. (iii) en notant que les $u_i\ (\underline{O}_{Z_i})$ sont bien quasi-coh.

\scanpage{66}
% archive sheet number: 65
% The entire draft is cancelled by a long diagonal stroke; two additional nearly vertical cancellation strokes flank the proof.
\deleted{Ceci donne une démonstration plus sympatique et élémentaire de 11.9.6. Quant à 11.9.3. (ii), il varésulter d'un}

\deleted{\uline{Lemme} (ou corollaire à la proposition précédente). Sous les hypothèses de la proposition, ne supposons plus néc que $A'$ soit projectif sur $A$, mais seulement ceci: pour tout quotient intègre $B=A/\underline{p}$ de $A$, posant $B'=B\otimes_A A'$, il existe un $h\quad B-(0)$ tel que $B_h$ soit \deleted{projectif}\inserted{libre} sur $A_h$. Alors la conclusion de la proposition reste encore valable.}

\deleted{Soit en effet $\handwritten{\phi}$ une section de $F'$ dont les images par les $u'$ sont nulles, prouvons $\handwritten{\phi}=0$, i.e. $\operatorname{supp}\handwritten{\phi}=\handwritten{\unsure{\varnothing|\phi}}$. Dans le cas contraire, il existerait un point maximal $x'$ dans $\operatorname{supp}\handwritten{\phi}$, correspondant à l'idéal premier $\underline{p}'$ de $A'$. Soit $\deleted{\unsure{\underline{p}=\underline{p}'\cap A|\underline{p}=\underline{p}'\,A}}\inserted{x}$ son image dans $X$, défini par l'idéal $\underline{p}$ trace de $\underline{p}'$ sur $A$, alors en vertu de l'hypothèse appliquée à $A/\underline{p}$, on voit que, quitte à remplacer $X$ par un voisinage ouvert affine de $x$, on peut supposer que $\unsure{\deleted{B}A|\deleted{A}B}'/\underline{p}\deleted{B}A'$ est \deleted{projectif}\inserted{libre} sur $A/\underline{p}$. Notons que nous pouvions choisir $x'$ de façon que $x$ soit également maximal dans $g(\operatorname{supp}\handwritten{\phi})$ (commencer par choisir $x$ maximal dans $g(\operatorname{supp}\handwritten{\phi})$, puis $x'$ maximal dans $g^{-1}(x)$), ce qui signifie que localisant en $\underline{p}$, le support de $\handwritten{\phi}_{\underline{p}}$ est contenu dans $\operatorname{Spec}(A'_{\underline{p}}/\underline{p}A'_{\underline{p}})$, donc $\handwritten{\phi}_{\underline{p}}$ est annulé par une puissance de $\underline{p}\ A'_{\underline{p}}$, ou encore que l'on peut trouver $h\quad A-\underline{p}$ tel que}

\scanpage{68}
% archive sheet number: 67
% Two long diagonal strokes cancel this draft. Unfilled sign slots are left as spacing.
\deleted{sous les conditions envisagées a) b) $\quad)$ $\quad)$, $B$ plat sur $A$ et à extension résiduelle finie,}
\[
\deleted{\text{\uuline{R}}(A)\quad\text{et}\quad\text{\uuline{R}}(B\otimes k)\ \text{géométrique}\quad\text{impliquent}\quad\text{\uuline{R}}(B).}
\]
\deleted{La condition supplémentaire qu'on vient de signaler pour \uuline{R} est satisfaite dans tous les exemples envisagés dans 7.5.3., à l'exception de (vii), où il est en défaut. En effet, on peut (même \deleted{sur un corps} pour une courbe $Y$ sur un corps algébriquement clos $k$ de caractéristique arbitraire) donner des exemples d'un morphisme étale $f\colon X\quad Y$, et d'un point $x$ de $X$ tel que $Y$ soit intègre en $y=f(x)$, mais $X$ ne le soit pas en $x$.\handwritten{/} Ce serait peut-être le moment de donner l'exemple (également signalé dans 6.15.7. et qui resert $X$ fois): on prend une courbe à point double ordinaire, et on recolle deux exemplaires du normalisé \unsure{d|de} de façon à avoir une courbe étale $X$ au dessus de $Y$. Un exemple analogue avec un anneau local \inserted{intègre} \uline{complet}, donc unibranche, mais en l'occurrence non géométriquement unibranche, est donné par $\inserted{A=}\text{\uuline{R}}\ X,Y\quad/(X^{\handwritten{2}}+Y^{\handwritten{2}})$. \deleted{Si $X$ est son spectre} Alors $A\otimes_{\text{\uuline{R}}}\underline{C}$ est étale sur $A$, mais non intègre. Ces exemples, qu'on semble rencontrer pour la première fois à propos de 6.15.7. pourraient peut-être remonter à cet endroit.\handwritten{/} Il est cependant plausible que \deleted{l'énoncé qui} la reponse à la question suivante soit affirmative :}

\deleted{Soit $A\quad B$ un homomorphisme local d'anneaux locaux noethériens, $B$ plat sur $A$, $B\otimes_A k\handwritten{/}k$ géométriquement intègre, $A$ \deleted{excellent (par exemple complet)} géométriquement unibranche (NB on peut se ramener au cas: $A$ normal). Si $A$ est de plus excellent (disons complet), alors $B$ est-il intègre ?}

\deleted{La reponse est affirmative quand le corps résiduel $k$ de $A$ est de car 0, comme on verra au N° 7.9. en utilisant la\deleted{\unsure{x}} résolution des singularités de Hironaka, ou lorsque $B$ est essentiellement de type fini sur $A$ (réf $\qquad$). Mais \inserted{dans le cas général,} même lorsque l'extension résiduelle est triviale, $A$\deleted{\unsure{s|x}} et $B$ complet, \deleted{\unsure{la technique dévis}} nous ignorons si la réponse est affirmative. On a cependant un résultat partiel ci-dessous.}

\scanpage{69}
% archive sheet number: 68
% This scan is handwritten throughout; the faint mirrored typescript is show-through from the reverse, not a second front-side text.
$\overline{\text{\unsure{un|Un}}}$ énoncé pour $X=\operatorname{Spec}\mathbb{Z}$, et dans le cas géométrique, pour $X=\operatorname{Spec}\mathbb{Z}/p\mathbb{Z}$. Dans ce dernier cas, on doit vérifier simplement la dualité entre
\[
H^i(\mathbb{F}_p,M)\quad\text{et}\quad H^{1-i}(\mathbb{F}_p,M^*)
\]
\[
(M^*=\unsure{\operatorname{Hom}|\mathscr{H}\mathrm{om}}(M,\unsure{\mathbb{Q}(0)|\mathbb{Q}\omega)}))
\]
à l'aide des ``\uline{cup-produits}'' et de l'isom.
\margin{\unsure{à développer un peu pour les cup-produits à cohomologie ``absolue''} $-\ -\ -$}
\[
H^1(\mathbb{F}_p,\mathbb{Q}(0))\simeq\mathbb{Q}
\]
et il suffit de faire la \unsure{vérification} $l$-adiquement, en appliquant le foncteur $T_\ell$. Cela revient : \unsure{la|le} \unsure{vérification suivante} [pour un \unsure{vectoriel|vecteur} de dim finie sur $\mathbb{Q}_\ell$ sur lequel $\pi=\mathbb{Z}$ opère]
\[
H^i(\pi,V)\quad\text{et}\quad H^{1-i}(\pi,V^*)\qquad\text{en dualité}
\]
\[
(V^*=\operatorname{Hom}_{\mathbb{Q}_\ell}(V,\mathbb{Q}_\ell)\ :\ \text{c'est une}
\]
trivialité.

\scanpage{70}
% archive sheet number: 69
% Several long oblique cancellation strokes cross the draft.
\deleted{\uline{Exemples} 7.5.3. \deleted{Les}\inserted{Chacune des} propriétés \inserted{\uuline{R}} suivantes satisfait aux condition a) et b) de 7.5.1., \deleted{pourvu que $B$ soit un} lorsque $B$ est un excellent anneau:}
\begin{enumerate}
\item[\deleted{(i)}] \deleted{$\operatorname{coprof}C\quad n$}
\item[\deleted{(ii)}] \deleted{$S_n$}
\item[\deleted{(iii)}] \deleted{Coh Mac}
\item[\deleted{(iv)}] \deleted{régulier}
\item[\deleted{(v)}] \deleted{normal $R_n$}
\item[\deleted{(vi)}] \deleted{normal}
\item[\deleted{(vii)}] \deleted{intègre}
\item[\deleted{(viii)}] \deleted{réduit}
\end{enumerate}

\deleted{En fait, pour (vii) et (viii) c'est vrai sans aucune hypothèse sur $B$, pour (i) (ii) (iii) il suffit de vérifier la condition (b)\inserted{et la cond univ cat dans (ii)}, \deleted{qui} ce \inserted{qui} est vérifiéé si \deleted{\unsure{tout que|tant que}} $B$ est quotient d'un régulier, ou satisfait la condition de Nagata sur le lieu singulier.\inserted{(et dans le cas (ii) est univ cat).} Cette condition est également suffisante pour faire marcher (iv) (v) (vi), pourvu que dans (v) et (vi) on suppose de plus $B$ universellement caténaire. Enfin, pour faire marcher (vi), on peut aussi prendre la condition ``$B$ est univ japonais et univ cat''.}
\margin{\deleted{gata}}

\deleted{\uline{Remarques} 7.5.4.\quad (i) comme 7.5.6. (i)}

\deleted{(ii) Comme 7.5.6. (ii) ; il s'agirait pratiquement des propriétés (i) \deleted{(ii)} (iii) (vii) ou (viii).}

\deleted{(iii) Dans la situation 7.5.1., on dispose le plus souvent d'une propriété supplémentaire de \uuline{R}, garantissant que si \deleted{$A\quad B$} $\text{\uuline{R}}(A)$ est vrai et les fibres de $\operatorname{Spec}(B)\quad\operatorname{Spec}(A)$ satisfont la condition \uuline{R} géométrique, \deleted{il en est de même} alors \deleted{($\operatorname{Spec}(B)$ étant} ($B$ étant plat sur $A$, rappelons-le) $\text{\uuline{R}}(B)$ est vrai. On peut donc tirer de 7.5.1. que}

\scanpage{81}
% archive sheet number: 80
% The typed draft is cancelled by oblique strokes across the page. Marginal pen additions are retained separately.
\deleted{et l'hypothèse de récurrence montre alors qu'on a l'inégalité}
\[
\deleted{\operatorname{long}M/tM\qquad\operatorname{long}M_{\underline{p}_i}\ ,}
\]
\deleted{l'égalité étant vraie sss les égalités analogues sont vraies pour $M'$ et $M''$, ce qui pqr l'hypothèse de récurrence signifie précisément que b) et c) sont vérifiés\inserted{pour $M$}. CQFD}

\deleted{\uline{Corollaire} 3. Sous les conditions de 1.\inserted{(ii)}, supposons que $\operatorname{long}(F/tF)_z$ soit égal à 1, alors il existe une seule composante irréductible $S_i$ de $S=\operatorname{supp}F$ passant par $Z$, et si $x$ est son point générique, on a $\operatorname{long}F_x=1$. De plus, $z$ est un point régulier de $\deleted{S_i}\inserted{\unsure{S_{\mathrm{réd}'}|S_{\mathrm{réd}},}}$ et $t$ engendre l'idéal maximal de $\underline{O}_{S_{\mathrm{réd}},z}$.}

\deleted{\uline{Proposition} 4\quad Sous les conditions de 1., soit $x\quad X$, supposons que $V(t)$ ne contienne aucune composante irréductible $S_i$ de $S=\operatorname{supp}F$ passant par $x$, et que $F/tF$ soit réduit en $x$. Alors $t$ est $F$-régulier en $x$, et $F$ est réduit en $x$.}

\deleted{Pour la première assertion, la démonstration est donnée dans 5.6.2. Le fait que $F_x$ n'ait pas de cycles premiers associés immergés passant par $x$ est \deleted{un cas} alors un cas particulier de \unsure{8|6}. ci-dessous. Reste à prouver que les long\inserted{u}eurs de $F$ aux points générisations maximales de $x$ dans $S$ sont égales à 1, ce qui est contenu dans le corollaire 3 précédent; on pourrait d'ailleurs ajouter que les générisations maximales de $x$ dans $\operatorname{supp}F/tF$ sont réguliers sur $S_{\mathrm{réd}}$ et que $t$ \unsure{engendre} l'idéal maximal des anneaux locaux de $S_{\mathrm{réd}}$ en ces points \ldots}

\deleted{\uline{Proposition} 5. Sous les conditions de 1., soit $T$ un cycle premier associé à $F$ (pas nécessairement maximal), $T'$ une composante de $T\quad V(t)$, \deleted{\unsure{supposons} de point générique $x$.} Supposons que $t$ est $F$-régulier au point générique $x$ de $T'$, alors \deleted{$T'$ est un} $x\quad\operatorname{Ass}F/tF$.}
\margin{\handwritten{$\left.\begin{gathered}\text{\unsure{Corollaire|Condition} pour}\\\unsure{x|t}\ \text{\unsure{soit|est}}\ F\ \text{\unsure{régulier}.}\end{gathered}\right]$}}

\deleted{La démonstration figure dans 8.1.4.}

\deleted{\uline{Corollaire} 6. Sous les conditions de 1., soit $x\quad X$, \deleted{et} supposons que $t$ soit $F$-régulier en $x$ et que $F/tF$ n'ait pas de cycles premiers associés immergés passant par $x$. Alors \inserted{\handwritten{$F$ n'a pas de cycles premiers immergés \unsure{passant} par $x$.}} les composantes irréductibles de $\operatorname{supp}F/tF$ passant par $x$ sont \unsure{e xactement|exactement} les composantes irréductibles passant par $X$ des $S_i\cap V(t)$, où les $S_i$ sont les comp irred de $\operatorname{supp}F=S$}

\scanpage{83}
% archive sheet number: 82
% Long oblique strokes cancel this draft.
\deleted{IV 7.5.}

\deleted{Je ne suis pas très content encore de la dernière version, car on n'énonce pas ce qu'on démontre vraiment, et par conséquent il y a des cas (tels la ``$\operatorname{copfof}\quad n$'', ou ``intègre'') qui échappent à l'énoncé tel qu'il est donné, mais sont inclus dans la démonstration. Aussi je propose de réénoncer les choses ainsi:}

\deleted{\uline{Théorème} 7.5.1.\quad Soit $A\quad B$ un homomorphisme local d'anneaux locaux noéthériens, $\text{\uuline{R}}(C)$ une propriété \deleted{d'anneaux locaux} faisant intervenir un argument qui est un anneau local \unsure{noehhérien|noethérien}. On suppose\deleted{\unsure{r|x}} que $(B,R)$ satisfait à :}

\deleted{a) Pour tout anneau \inserted{local} $B^{\deleted{\unsure{'|x}}}$ localise d'une $B$-algèbre finie, et tout élément régulier $t\quad\underline{r}(C)$, $\text{\uuline{R}}(C/tC)$ implique $\text{\uuline{R}}(C)$.}

\deleted{b) Si $C$ est une $B$-algèbre finie, alors $U_{\text{\uuline{R}}}(\operatorname{Spec}(C))$ est ouvert.}

\deleted{On suppose de plus que $A$ satisfait les conditions suivantes\deleted{\unsure{x}} (remplies si $A$ est ``excellent'' (7.8.2.) en particulier si $A$ est complet):}

\deleted{\handwritten{$\alpha$}) $A$ est universellement japonais.}

\deleted{\handwritten{$\beta$}) Pour toute algèbre finie intègre $C$ sur $A$, $\operatorname{Spec}(C)$ contient un ouvert non vide qui est régulier,}

\deleted{Soit $k$ le corps résiduel de $A$, \inserted{et supposons le corps résiduel de $B$ fini sur celui de $A$.} Supposons $B$ plat sur $A$ et que $B\otimes_A k$ sur $k$ satisfait la condition ``$Q$ géométrique''; Alors toutes les fibres de $\operatorname{Spec}(B)$ sur $\operatorname{Spec}(A)$ satisfont la condition $Q$ géométrique.}

\deleted{Démonstration. Les hypothèses faites sur $A\quad B$ sont stables par changement de base fini $A\quad A'$, pourvu qu'on prenne soin de reformuler les hypothèses pour $A,B$ semi-locaux, l'homomorphisme $A\quad B$ semi-local (i.e. l'image inverse d'un idéal maximal est un idéal maximal) et qu'on remplace le corps résiduel de $A$ par les corps résiduels de $A$. \deleted{(NB alors} (NB cette extension ne mérite pas de figurer dans l'énoncé du théorème, c'est juste un intermédiaire technique pour la démonstration). On est ainsi ramené au cas $A$ intègre et à voir que la fibre générique satisfait}

\scanpage{85}
% archive sheet number: 84
% Three large diagonal cancellation strokes cross the typescript.
\deleted{\uline{Remarques} \inserted{\handwritten{4}} a) On se gardera de croire que lorsque $A$ est un anneau \unsure{noethérien}, $B$ une $A$-algèbre plate de type fini, $(J_i)_{i\quad I}$ une famille d'idéaux de $A$ dont l'intersection est nulle, alors l'intersection des $J_iB$ soit nulle. (On a déjà signalé l'exemple d'un anneau de valuation discrète $A$, muni des puissances \inserted{$\mathbb{J}_i$} de son idéal maximal, et $B$ étant le corps des fractions). \deleted{Le fait} En effet, l'hypothèse que $\quad J_i=0$ n'implique nullement que $\quad J_i=0$ !}

\deleted{b) Soit $X$ un préschéma sur un corps $k$, et soit $u_i\colon Z_i\quad X$ une famille schématiquement dominante. Nous ignorons si pour tout changement de base $S\quad\operatorname{Spec}(k)$, la famille des $u_{i\,S}\colon\unsure{Z_{i\,S}|Z_{A\,S}}\quad X_S$ est sch. dominante\inserted{cf définition 5}. On peut \deleted{se ram} montrer (par des arguments du type de ceux utilisés ci-dessous) qu'il suffit pour cela que cette propriété soit satisfaite chaque fois que $S$ est le spectre d'un corps $k'$ extension finie séparable de $k$. D'autre part, nous verrons ci-dessous que la réponse est affirmative lorsque $X$ est localement noethérien.}

\deleted{\uline{Définition} 5. Sous les notations de 1., supposons que $X$ soit un $S$-préschéma. On dit que $(u_i)$ est \deleted{relativement} schématiquement dominante relativement à $S$ \deleted{(ou que la famille des sous-préschémas $Z_i$, s'il s'agit de tels, est sch dense dans $X$ rel à $S$)} si pour tout $S'\quad S$, la famille des $u_{i\,S'}\colon Z_{i\,S'}\quad X_{S'}$ est sch dominante \deleted{(resp ou que la famille} Quand les $Z_i$ sont des sous-préschémas de $X$ \ldots, on dit aussi que la famille \deleted{des} $\unsure{(Z_i)|\{Z_i)}$ est sch dense dans $X$ rel à $S$.}

\deleted{NB C'est une notion locale sur $X$ et sur $S$.}

\deleted{\uline{Lemme} 6\quad Avec les notations précédentes, supposons $X$ et les $Z_i$ plats sur $S$, soit $S_{\unsure{o|0}}$ un sous-préschéma de $S$ défini par un Idéal Nilpotent. Si la famille des $u_{i\,S_o}$ est sch dominante, il en est de même de la famille des $u_i$. On peut resper avec ``rel à $S$''.}

\deleted{La deuxième assertion résulte trivialement de la première et des définitions. Pour la première, cf SGAD IX 4.2.}

\scanpage{97}
% archive sheet number: 96
\begin{center}
Filtration et poids des espaces de cohomologie des variétés algébriques

par A.GROTHENDIECK
\end{center}

0. \uline{Introduction}. On a remarqué, depuis \deleted{longtemps}\inserted{l'invention de la cohomologie $\handwritten{\ell}$-adique \handwritten{\unsure{([\quad]|[\quad])}}}, que\deleted{\unsure{,|x}} pour une variété algébrique $X$ définie sur un corps $k$, de clôture algébrique $\bar{k}$, les opérations naturelles du groupe de Galois $\deleted{\unsure{\operatorname{Gal}k|G_k|G(k)}}\handwritten{\pi}=\operatorname{Gal}(\bar{k}/k)$ sur la cohomologie $\handwritten{\ell}$-adique $H(X_{\bar{k}},\underline{Z}_{\handwritten{\ell}})$ de la variété ``géométrique'' $X_{\bar{k}}$ définissait sur cette \deleted{dernière}\inserted{\handwritten{cohomologie}} une structure supplémentaire remarquable, donnant lieu notamment aux remarquables conjectures de Tate \handwritten{[\quad]}. Il semble probable que ces dernières soient n\deleted{\unsure{o|e}}\handwritten{e}t\deleted{\unsure{ablem|able}}\handwritten{tem}ent plus éloignées d\handwritten{e}\deleted{leur solution} leur solution que les conjectures de Weil \handwritten{[\quad,\quad]}, qui, comme on sait, peuvent s'interpréter également en termes d'actions galoisiennes sur des groupes de cohomologie $\handwritten{\ell}$-adiques. Notre propos dans le présent travail est de montrer comment on peut, des conjectures de Weil et (le cas échéant) de la résolution des singularités, tirer des conséquences de nature ``géométrique'' en apparence, i.e. relatives à un corps de base algébriquement clos. Le principe de cette méthode\inserted{\deleted{(qui remonte à 1964)}}, qui consiste donc essentiellement à ramener un énoncé géométrique au cas d'un corps de base fini, est expliqué au N°1, et des exemples d'application de cette méthode sont donnés dans les numéros 2 à 7. \deleted{\unsure{Le pr}} L'idée de ces raisonnements, liés à\deleted{\unsure{ux|la}}\inserted{la} considération\deleted{s}\inserted{systématique} des ``poids'' des modules galoisiens provenant de la cohomologie $\handwritten{\ell}$-adique, remonte à 1964; des exemples \deleted{du même principe}\inserted{\handwritten{analogues}} se trouvent développés dans SGA 7\handwritten{(}I et II\handwritten{)}; dans le séminaire oral\deleted{e} était\inserted{\handwritten{développé}} également l'\unsure{exempbe|exemple} du n°2 du présent travail.

Bien entendu, la portée des ``résultats'' de notre travail est nécessairement \deleted{limitée}\inserted{\handwritten{sujette à caution}}, puisqu'ils dépendent de conjectures non prouvées à l'heure actuelle, savoir les conjectures de Weil et la résolution des singularités (sous la forme forte de Hironaka) en caractéristique

\scanpage{98}
% archive sheet number: 97
quelconque. \inserted{Par expérience personnelle,} Nous pensons cependant\deleted{,} que la compréhension\deleted{, heuristique quoique pour le moment largement heuristique,} obtenue grâce aux arguments de poids et de niveaux, quoiqu'elle soit pour l'instant entièrement heuristique, \deleted{\unsure{est et}} peut être la source de la découverte de propriétés importantes et inattendues pour la cohomologie des variétés algébriques, \deleted{même si} et qu'elle mérite pour cette raison de ne pas être dédaignée. De plus, dans le cas des corps de caractéristique nulle, \deleted{des travaux récents fort \unsure{remar|rem} \unsure{les|las} magnifiques t} le récent et magnifique travail de P. DELIGNE sur la théorie de Hodge\inserted{généralisée} des variétés algébriques quelconques sur le corps $\underline{C}$ des nombres complexes\deleted{,} justifie\inserted{dès maintenant} par voie transcendante le ``yoga'' des poids et niveaux dans le cadre de la cohomologie de De Rham\handwritten{,} (tout en ouvrant à la théorie cohomoligue des variétés algébriques\inserted{en caractéristique nulle} des voies entièrement nouvelles). Il \deleted{est}\inserted{semble} probable d'autre part que pour le cas des \deleted{Pour terminer cette introduction} variétés algé\deleted{\unsure{bu|br}}\handwritten{br}iques sur les corps de caractéristique non nulle, une justification de ce même yoga doive attendre la \deleted{\unsure{résolution|resolution}}\inserted{preuve} des conjectures de Weil (ou mieux, des ``conjectures standart'' de \handwritten{[\quad]}).

Pour \deleted{signaler} terminer cette introduction, signalons que les filtrations par poids et niveaux des espaces de cohomologie $\handwritten{\ell}$-adiques, pour $\handwritten{\ell}$ variable, sont certainement relié\deleted{\unsure{s}}\handwritten{es} entre elles de la façon habituelle, \deleted{donnant}\inserted{et doivent donner} en particulier des espaces vectoriels sur $\underline{Q}_{\handwritten{\ell}}$ \handwritten{d}ont le \deleted{\unsure{ran|rang}} rang est indépendant de $\handwritten{\ell}$. En termes de la théorie conjecturale des motifs (suspendue entre autres aux ``conjectures standart'' déjà mentionnées), ceci devrait s'exprimer par le fait que les structures de poids et de niveau\deleted{se placent}\inserted{des cohomologies $\handwritten{\ell}$-adi\handwritten{ques} se placent déjà} au niveau de\deleted{\unsure{s motifs|motifs} déjà} la \uline{cohomologie} \deleted{des} \uline{motivique}, \deleted{ce qui dans la}\inserted{Cet\deleted{\unsure{te|t}} point de vue est en tous cas en accord avec la} théorie de Deligne, \deleted{\unsure{car}} puisque sa ``filtration par le poids'' de la cohomologie de De Rham provient d'une filtration de la cohomologie à coefficients \uline{rationnels}\handwritten{.}\deleted{(et non seulement}

\scanpage{100}
% archive sheet number: 99
\uline{2. Formalism of Algebraic Cohomology Classes}.

\uline{Cohomology Class Associated to a Cycle}. Let $C^*(X)$ be the intersection\inserted{\handwritten{ring}} of cycles on $X$,\inserted{\handwritten{modulo rational equivalence \unsure{(the Chow ring of $(X)$|(the Chow ring of $(K)$}}} tensored by $Q$. Then, for each $i\handwritten{\leq}n$, we define a homomorphism
\[
\handwritten{\rho_X:C^i(X)\longrightarrow H_\ell^{2i}(X)}
\]
as follows. Let $Z$ be an irreducible subscheme of codimension $i$. The l-adic Poincare duality theorem asserts that the cup-product pairing
\[
\handwritten{H_\ell^j(X)\times H_\ell^{2n-j}(X)\longrightarrow H_\ell^{2n}(X)\overset{\unsure{\kappa_X|K_X}}{\overset{\sim}{\longrightarrow}}Q_\ell\quad(j=0,\ldots,2n)}
\]
is a perfect duality of finite dimensional vector spaces, and so the element $\handwritten{\rho_X(Z)\in H_\ell^{2i}(X)}$ is uniquely determined by the condition
\[
\handwritten{\kappa_X\bigl(\alpha\cdot\rho_X(Z)\bigr)=\kappa_Z\bigl(\operatorname{res}_Z(\alpha)\bigr)\quad\text{for all }\alpha\in H_\ell^{2n-2i}(X).}
\]

We then extend $\handwritten{\rho_X}$ to all of $C^{\handwritten{i}}(X)$ by linearity. Further, this homomorphism takes intersection product into cup-product, and thus it extends to a homomorphism of \handwritten{``}algebras\handwritten{''} \deleted{on}\inserted{\handwritten{over}} $Q$\unsure{.}
\[
\handwritten{\rho_X:C^*(X)\longrightarrow H_\ell^*(X)..}
\]

The image of a cycle under this homomorphism .is called the \uline{cohomology class} associated with \unsure{thet|that} cycle, and we denote $\handwritten{\rho_X}(C^*(X))$ by $C_{\handwritten{\ell}}^*(X)$. A cycle is defined to be $Q_{\handwritten{\ell}}$-cohomologically equivalent to $0$ if it is in the kernel of the homomorphism $\handwritten{\rho_X.}$
